% ECONOMICS_PROBLEM: {"id": "C23-406", "title": "Staggered adoption with dynamic heterogeneous effects", "jel_code": "C23", "source_id": "OP-0107"}
% !TeX program = xelatex
\documentclass[11pt,a4paper]{article}
\usepackage[margin=23mm]{geometry}
\usepackage{fontspec}
\setmainfont{Latin Modern Roman}
\usepackage{amsmath,amssymb,mathtools}
\usepackage{xcolor,enumitem,booktabs,longtable,array,xurl,hyperref}
\definecolor{muted}{HTML}{53636C}
\hypersetup{colorlinks=true,urlcolor=blue,linkcolor=blue}
\setlength{\parindent}{0pt}
\setlength{\parskip}{5pt}
\newcommand{\R}{\mathbb R}
\newcommand{\E}{\mathbb E}
\newcommand{\N}{\mathbb N}
\newcommand{\ind}{\mathbf 1}
\newcommand{\Var}{\operatorname{Var}}
\newcommand{\Cov}{\operatorname{Cov}}
\newcommand{\supp}{\operatorname{supp}}
\DeclareMathOperator*{\argmin}{arg\,min}
\DeclareMathOperator*{\argmax}{arg\,max}
\newcommand{\entrymeta}[3]{{\sffamily\small\color{muted}\textbf{\texttt{#1}}\quad|\quad #2\par #3\par}}
\newcommand{\Problemref}[1]{\texttt{#1}}
\newcommand{\JELref}[2]{\texttt{#2} (JEL \texttt{#1})}
\begin{document}
\section*{Staggered adoption with dynamic heterogeneous effects}
\textbf{Problem ID:} \texttt{C23-406}\quad \textbf{JEL:} \texttt{C23}\par
% BEGIN ECONOMICS STATEMENT
\label{problem:OP-0107}
\label{atlas:prob:E7}

\noindent\textbf{Original identifier:} E7 (atlas item 107). \textbf{Field:} Econometrics and causal inference. \textbf{Classification:} editorial research family with established restricted solutions; current open status not certified.

\subsection*{Definitions and assumptions}

In an independent-unit panel $(Y_{it},D_{it},X_i)$, $t=1,\ldots,T$, define adoption cohort $G_i\in\{2,\ldots,T,\infty\}$ and initially assume absorbing treatment $D_{it}=\mathbf1\{t\ge G_i\}$. Let $Y_{it}(g)$ be the potential outcome under adoption at $g$, $Y_{it}(\infty)$ under no treatment, and impose no anticipation: $Y_{it}(g)=Y_{it}(\infty)$ for $t<g$. Define $\operatorname{ATT}(g,t)=\mathbb E[Y_{it}(g)-Y_{it}(\infty)\mid G_i=g]$ for $t\ge g$. Comparison groups must be untreated and have covariate overlap; conditional parallel trends equates untreated changes for cohort $g$ and its stated eligible controls.

\subsection*{Formal research question}

For prespecified nonnegative weights $\omega_{g,t}$ summing to one over observed cohort-time cells, identify and infer $\theta=\sum_{g,t}\omega_{g,t}\operatorname{ATT}(g,t)$. This no-anticipation, absorbing-adoption baseline is addressed by existing methods. For the extension, fix an integer anticipation horizon $a\in\{0,\ldots,T-1\}$ and specify full announced treatment and exposure regimes $d=(d_1,\ldots,d_T)$ and $h=(h_1,\ldots,h_T)$, allowing reversals. Potential outcomes $Y_{it}(d,h)$ may depend on the regimes through date $\min(t+a,T)$; imposing dependence only through $t$ would rule out the permitted anticipation. Define a fixed-population average contrast $\mathbb E[Y_{it}(d,h)-Y_{it}(d',h')]$ between two feasible announced regimes. Characterize support, information and sequential comparison restrictions sufficient for identification; with supplied nonnegative bounds $\delta_{g,t}$ on untreated-trend violations, derive a sharp sensitivity set and confidence regions with uniform $1-\alpha$ coverage, $\alpha\in(0,1)$, over a declared observed-law and regime-response class.

\subsection*{Interpretation and scope}

Weights determine the policy question: equal cohort weights, population weights and event-time weights generally give different answers. Estimated weights require uncertainty accounting. Comparisons to already treated units may incorporate treatment dynamics, so a two-way-fixed-effects coefficient need not equal the desired weighted ATT; decomposition weights and implicit weights on cohort-time effects must be distinguished. Event-study pretrend tests cannot prove future parallel trends and can themselves be contaminated by heterogeneous effects. Treatment reversals require history-dependent outcomes and fresh identifying conditions. Spillovers remove untreated controls unless exposure, rather than own treatment only, defines eligibility.

\subsection*{Source audit and status}

[S1] decomposes staggered-treatment two-way-fixed-effects estimators; [S2] provides group-time identification, estimation and inference; [S3] handles heterogeneous dynamic effects in event studies. Thus the original aggregation question has major established solutions under absorbing adoption and parallel trends. The history, anticipation and spillover extension is an editorial research family. None of these 2021 citations certifies that a particular unrestricted extension is still open in October 2026.

\subsection*{References}

\begin{enumerate}[label={[\textnormal{S}\arabic*]},leftmargin=*]

\item Andrew Goodman-Bacon (2021). \emph{Difference-in-differences with variation in treatment timing}. Journal of Econometrics 225(2), 254–277. \url{https://doi.org/10.1016/j.jeconom.2021.03.014}. \textbf{Locator:} Publisher or institutional abstract and bibliographic record. \textbf{Support and limits:} Decomposes two-way-fixed-effects comparisons; already diagnoses a central flaw.

\item Brantly Callaway and Pedro H. C. Sant’Anna (2021). \emph{Difference-in-Differences with multiple time periods}. Journal of Econometrics 225(2), 200–230. \url{https://doi.org/10.1016/j.jeconom.2020.12.001}. \textbf{Locator:} Publisher or institutional abstract and bibliographic record. \textbf{Support and limits:} Identification, estimation and inference for group-time effects under parallel trends and other restrictions.

\item Liyang Sun and Sarah Abraham (2021). \emph{Estimating dynamic treatment effects in event studies with heterogeneous treatment effects}. Journal of Econometrics 225(2), 175–199. \url{https://doi.org/10.1016/j.jeconom.2020.09.006}. \textbf{Locator:} Publisher or institutional abstract and bibliographic record. \textbf{Support and limits:} Interaction-weighted event-study results address heterogeneous absorbing adoption; not all reversals or interference.

\end{enumerate}

\subsection*{Common conventions: Source mathematical conventions}

Unless an entry states otherwise, agent and firm sets are finite. State and action spaces are measurable, strategies and outcome maps are measurable, probability laws are specified, and expectations exist for the targets under discussion. Infinite-horizon discount factors lie in $(0,1)$ and discounted objectives are finite. Norms, tolerances and complexity input sizes must be specified for computational comparisons. Constants or primitives that appear locally are inputs, not empirical facts asserted by this catalogue.

\textbf{Identification.} A statistical model $m\in\mathcal M$ implies an observable law $P_m$ and target $\theta(m)$. At law $P$, its identified set is
\[
\Theta(P;\mathcal M)=\{\theta(m):m\in\mathcal M,\ P_m=P\}.
\]
Point identification holds when this set is a singleton, or equivalently when $P_{m_1}=P_{m_2}$ implies $\theta(m_1)=\theta(m_2)$ for all compatible models. Sharp bounds mean bounds attained, or approached when the boundary is not attained, by compatible models. Writing this set is a definition, not a solution: the substantive task is to establish informative restrictions and derive or estimate the set. An unrestricted-confounding model may give only trivial bounds.

\textbf{Inference.} Identification, estimability and valid uncertainty quantification are separate questions. Fix $\alpha\in(0,1)$. A confidence procedure $C_n$ over a declared law class $\mathcal P$ has asymptotic uniform coverage at level $1-\alpha$ when
\[
\liminf_{n\to\infty}\inf_{P\in\mathcal P}\Pr_P\{\theta(P)\in C_n\}\geq1-\alpha.
\]
For set-identified targets the coverage event must be defined separately, for example coverage of the full identified set. If an entry uses identification-indexed models instead of uniquely indexed laws, the same coverage convention is applied over those models. Pointwise limits do not establish uniform validity, and asymptotic validity does not guarantee finite-sample precision. Sample sizes, dependence structures and weak-identification sequences are part of each application.

\textbf{Counterfactuals.} $Y_i(a)$ is a potential outcome under a specified intervention, including its timing, implementation and versions. It is not $Y_i$ conditional on voluntarily choosing $a$. A full intervention vector permits interference; shorthand scalar treatment notation does not silently impose no interference. Assignment, selection, attrition, anticipation and measurement must be specified. A claim about an instrument's local effect is not automatically a population policy effect.

\textbf{Economic equilibrium.} A counterfactual model includes best responses, constraints, feasibility, market clearing, state transitions and expectations consistent with the stated information. When several equilibria exist, either a declared selector is an input or the target is set-valued. Comparisons across policy regimes must use the stated selection convention. Reduced-form relationships are not presumed invariant after an equilibrium-changing intervention.

\textbf{Optimization and welfare.} Optimization is over the stated feasible set. If attainment is not established, $\sup$ or $\inf$ and $\varepsilon$-optimal choices replace an asserted maximizer or minimizer. For example, with value $V=\sup_{a\in\mathcal A}W(a)<\infty$, an $\varepsilon$-optimal set is $\{a\in\mathcal A:W(a)\geq V-\varepsilon\}$ for $\varepsilon>0$. Benefits, costs, risk and health or environmental outcomes can be aggregated only using declared units and conversion weights. Interpersonal utility weights, the treatment of fiscal transfers and foreign profits, discounting, and the behavioral welfare criterion are explicit inputs. A comparative-static derivative additionally requires existence and appropriate differentiability of the selected counterfactual.

\textbf{Sources.} Local labels [S1], [S2], and so forth restart in every question. Locators identify the relevant abstract, model, section or result at the level actually checked; a bibliographic or abstract-level verification is not represented as inspection of an entire proof. Working papers are distinguished from later publications and changing versions. Source summaries are paraphrased. All URL checks and editorial formulations refer to the audit date stated above.
% END ECONOMICS STATEMENT
\end{document}
